% ECONOMICS_PROBLEM: {"id": "G14-317", "title": "High-frequency trading and market design", "jel_code": "G14", "source_id": "OP-0099"}
% FIRST_POSTED: 2026-10-09
% ATTEMPT_STATUS: unresolved
% ENTRY_KIND: research_attempt
\documentclass[11pt]{article}
\usepackage[margin=1in]{geometry}
\usepackage{amsmath,amssymb,amsthm,hyperref}
\newtheorem{proposition}{Proposition}
\newtheorem{lemma}{Lemma}
\newcommand{\E}{\mathbb E}
\newcommand{\R}{\mathbb R}
\newcommand{\Var}{\operatorname{Var}}
\newcommand{\Cov}{\operatorname{Cov}}
\begin{document}
\section*{High-frequency trading and market design}
\textbf{Research attempt by GPT-6 (Codex), 9 October 2026.}
The procedure was to restate the canonical target, inspect its cited primary source,
derive a tractable result, and test the assumptions and remaining gap.
These calculations are not a claim of novelty or independent proof verification.

\subsection*{Target, conventions and source context}
A market design $d$ controls messages, speed priority, clearing and information.
The target is to maximize total expected surplus net of real costs, subject
to price-error and execution-failure bounds, allowing strategic entry and
venue migration. Price sampling and the efficient-value benchmark must be
fixed before comparing designs. Budish, Cramton and Shim (2015), inspected
in its publisher full text, analyzes speed races and frequent batch auctions
in a particular model. Its results motivate, but do not prove, optimality
for the larger design menu in the canonical statement. The calculations
below are explicit restricted benchmarks; no observed exchange effect is
estimated here.

\subsection*{A solved speed-investment contest}
There are $n\geq2$ risk-neutral trading firms. Firm $i$ purchases speed
resource $e_i\geq0$ at unit cost. A fixed rent prize $Q>0$ is allocated
with probability $e_i/\sum_j e_j$ when total effort is positive. Payoff is
$Qe_i/\sum_j e_j-e_i$. At zero total effort split the prize equally.
The prize is a transfer from other market participants, not new output.

\begin{proposition}
The symmetric positive-effort equilibrium has
\[
 e_i^*=Q\frac{n-1}{n^2},\qquad
 \sum_i e_i^*=Q\frac{n-1}{n}.
\]
Counting the prize as a social benefit in addition to the payer's loss
would double-count it. The speed resources are a real loss in this model.
\end{proposition}
\begin{proof}
Given other firms' total effort $E_{-i}>0$, payoff has derivative
$QE_{-i}/(e_i+E_{-i})^2-1$ and strictly negative second derivative.
Substitute $E_{-i}=(n-1)e$ into the first-order condition to obtain the
stated $e$. Concavity proves global optimality, including comparison with
zero effort. Each player's net payoff is $Q/n^2>0$. Summing receipts
and the counterpart payer's loss cancels $Q$ exactly; resource expenditures
remain. We assert the displayed symmetric equilibrium, not an unproved
uniqueness theorem for arbitrary contest variations.
\end{proof}
Free entry would additionally compare $Q/n^2$ with fixed entry costs and
integer entry conditions. Simply holding $n$ fixed while claiming to
identify market-wide adaptation would be inconsistent.

\subsection*{Batch spacing with two binding design constraints}
Consider a separate batch benchmark indexed by interval $\Delta>0$.
Efficient value follows a Brownian motion with variance rate $s^2>0$.
At every clearing the published price equals the current efficient value;
between clearings the sampled quote is that last price. Sampling time is
uniform over a cycle and independent of value innovations. Thus
\[
 V(\Delta)=\frac1\Delta\int_0^\Delta s^2t\,dt
 =\frac{s^2\Delta}{2}.
\]
This describes stale quotes; sampling transaction prices only at clearing
would give zero error in this particular frictionless information model.
The distinction prevents silently changing the accuracy target.

Investor orders arrive at a stationary rate $\Lambda$ and each unit of
waiting time incurs real cost $c$. Mean waiting is $\Delta/2$, so the
waiting-cost rate is $A\Delta$ with $A=\Lambda c/2$. Each auction costs
$K>0$ real resources, giving rate $K/\Delta$. Suppose a required processing
task has an independent exponential service time with rate $\mu$; define
failure as service taking longer than the allotted interval. Then
$\Pr(A(\Delta))=e^{-\mu\Delta}$. This is an explicit technical failure
model, not a claim that market crashes have exponential probabilities.

For accuracy tolerance $\bar V>0$ and failure tolerance
$0<\bar p<1$, feasible intervals satisfy
\[
 L=\frac{\log(1/\bar p)}\mu\leq\Delta\leq
 U=\frac{2\bar V}{s^2}.
\]
If $L>U$ the design menu is infeasible. Otherwise maximizing net surplus
is equivalent to minimizing $A\Delta+K/\Delta$ on $[L,U]$.
For $A>0$ the unique optimum is
\[
 \Delta^*=\min\{U,\max\{L,\sqrt{K/A}\}\}.
\]
The derivative is $A-K/\Delta^2$ and the second derivative is
$2K/\Delta^3>0$, proving the claim and all boundary cases. If waiting is
costless ($A=0$), $U$ minimizes operating costs. Unbounded feasible
intervals and zero resource cost need separate attainment checks.

For uncertain $(A,K,s^2,\mu)$ in a rectangle with positive finite bounds,
the robust objective uses $\overline A,\overline K$ and robust constraints
use $\overline{s^2},\underline\mu$. This follows from coordinatewise
monotonicity, not a generic minimax interchange. A correlated uncertainty
set need not have all these adverse endpoints jointly attainable and
requires its actual support function.

\subsection*{An identification failure from venue reallocation}
Let fraction $q$ of orders route to the redesigned venue. Write observed
average execution quality as
$\overline L=qL_1+(1-q)L_0$, where each $L_j$ is conditional on the
population routed to venue $j$. A change in $\overline L$ decomposes into
within-venue changes and composition terms. For a simple explicit example,
two order types have fixed losses 0 and 2 at either venue. Before a rule
change the treated venue receives both types equally and has mean loss 1;
afterwards it receives only the easy type and has mean loss 0. Its design
has changed no type's outcome. A naive comparison nevertheless attributes
an improvement of 1 to the rule. Reversing selection generates the opposite
conclusion without changing within-type technology.

Randomizing rules across sufficiently isolated markets can identify their
equilibrium outcomes if each market is a treatment unit, assignment is
exogenous, and cross-market liquidity-provider spillovers are excluded or
modeled. Randomizing orders within a common venue identifies a different
response because the shared equilibrium and routing incentives persist.
These are substantive restrictions, not corrections a larger sample alone
can supply.

\subsection*{Verification and remaining gap}
The contest establishes resource dissipation under its payoff rule; the
batch optimization establishes an explicit feasible spacing. They are
separate modules: we have not proved that the batch mechanism eliminates
the contest in every equilibrium of the same underlying investor economy.
That missing link matters. A complete attempt must jointly specify order
strategies, information, entry, routing and clearing, then derive $W,V$
and failure probabilities from that equilibrium. The strongest next step
is a two-venue model with finite order types and endogenous entry, to test
whether an equilibrium shifts speed expenditure outside the regulated
venue. \textbf{Final status: UNRESOLVED.}

\section{Completion Estimate}
27\%

This is a post hoc, heuristic estimate for the full catalogue target.
The attempt proves speed-resource dissipation in a contest and solves batch spacing under explicit price-error and execution-failure constraints, including rectangular robustness. It does not join these modules into a common equilibrium with strategic orders, information, entry, routing, venue migration, and causally identified design effects.

\subsection*{Source}
E. Budish, P. Cramton and J. Shim (2015), \emph{The High-Frequency Trading
Arms Race: Frequent Batch Auctions as a Market Design Response},
Quarterly Journal of Economics 130(4), 1547--1621; Sections VI--VIII.
\url{https://academic.oup.com/qje/article/130/4/1547/1916146}.

\end{document}
